\documentclass[a4paper, 12pt]{article}
\usepackage[margin=1in]{geometry}
\usepackage[utf8]{inputenc}

\usepackage{amsmath, amssymb, amsfonts, bm, cancel, mathrsfs, mathtools, esint, bbm, tcolorbox,xcolor, vntex}

\numberwithin{equation}{section} 
\setlength{\parindent}{0pt}
\setlength{\parskip}{1em} 

\title{\textbf{Một số bài toán sưu tầm}}
\author{Đặng An Sang}
\date{}

\newtcolorbox{BoxRed}{colback=red!10, colframe=red!60, boxrule=0.4mm}

\begin{document}
	\maketitle
    
	\begin{BoxRed}
	\textbf{\textcolor{red}{Bài toán 1.}}
		Đặt $\ x_n = \sum\limits_{k = 1}^{n} {k \sin \dfrac{k}{n^3}}, \ \forall n \in \mathbb{N}^* \quad $
		\\Chứng minh rằng dãy $\{ x_n \}_{n=1}^{+\infty}$ có giới hạn hữu hạn và tìm giới hạn đó.\\
	\end{BoxRed}
	\textbf{\textcolor{red}{LỜI GIẢI}} \\
	Trước hết, ta chứng minh bổ đề: $\sin x \leq x, \ \forall x \in [0, +\infty)$\\
	Thật vậy, xét hàm số $f(x) = x - \sin x, \ \forall x \in [0, +\infty)$ ta có:\\
	$f'(x) = 1 - \cos x \geq 0, \ \forall x \in [0, +\infty)$ \\
	Từ đó, ta suy ra $min f(x) = 0$ khi $x = 0$ và $f$ đơn điệu tăng trên $[0, +\infty],$ qua đó \\$ \ f(x) = x - \sin x \geq 0, \ \forall x \in [0, +\infty)$
	\\Vậy, bổ đề được chứng minh.
	
	Quay trở lại bài toán, áp dụng bổ đề trên ta có:
	\[0 \leq \left| x_n - \dfrac{1}{3} \right| = \left| \sum_{k=1}^{n} k \sin \dfrac{k}{n^3} - \dfrac{1}{3} \right| 
	\leq \left| \sum_{k=1}^{n} \dfrac{k^2}{n^3} - \dfrac{1}{3} \right| 
	= \left| \dfrac{n(n+1)(2n+1)}{6n^3} - \dfrac{1}{3} \right|\]
	\\Mặt khác, ta luôn có 
	\[\lim \left| \dfrac{n(n+1)(2n+1)}{6n^3} - \dfrac{1}{3} \right|
	= \lim \left| \dfrac{\left(1 + \dfrac{1}{n}\right)\left(2 + \dfrac{1}{n}\right)}{6} - \dfrac{1}{3} \right| = 0 = \lim 0 \]
	Do đó, theo BĐT trên và nguyên lý kẹp suy ra $\lim x_n = \dfrac{1}{3}$
    \newpage
    \begin{BoxRed}
	\textbf{\textcolor{red}{Bài toán 2.}}
		Cho dãy $\{ x_n \}_{n=1}^{+\infty}$ được xác định bởi $$x_1 = 1, \ x_{n+1} = \sqrt{x_n(x_n + 1)(x_n + 2)(x_n + 3) + 1}, \ \forall n \in \mathbb{N}^*$$
		Đặt $y_n = \sum\limits_{i=1}^n \dfrac{1}{x_i + 2}, \ \forall n \in \mathbb{N}^*$. Tính $\lim y_n$
	\end{BoxRed}
	\textbf{\textcolor{red}{LỜI GIẢI}} \\
    Dễ thấy $u_n > 0, \ \forall n \in \mathbb{N}^*$. Từ công thức truy hồi ta có\\ 
	$u_{n+1} = \sqrt{(u_n^2 + 3u_n + 1)^2} = u_n^2 + 3u_n + 1$\\
	$\implies u_{n+1} - u_n = (u_n + 1)^2 > 0$\\
	$\implies \{u_n\}$ là một dãy tăng. Giả sử rằng bị chặn trên và hội tụ về $L > 0$. Cho qua giới hạn hai vế từ công thức truy hồi ta có\\
	$ L = L^2 + 3L + 1 \implies L = -1$ (mâu thuẫn)\\
	$\implies \lim u_n = +\infty.$ \\
	Cũng từ công thức truy hồi ta có\\ 
	$u_{n+1} + 1 = (u_n + 1)(u_n + 2) \implies \dfrac{1}{u_{n+1} + 1} = \dfrac{1}{u_n + 1} - \dfrac{1}{u_n + 2}$
	
	$\implies y_n = \sum\limits_{i=1}^n \dfrac{1}{x_i + 2} = \sum\limits_{i=1}^n \left(\dfrac{1}{x_i + 1} - \dfrac{1}{x_{i+1} + 1} \right) = \dfrac{1}{x_1 + 1} - \dfrac{1}{x_{n+1} + 1}$
	
	$\implies \lim y_n = \dfrac{1}{2}$
    \newpage
    \begin{BoxRed}
	\textbf{\textcolor{red}{Bài toán 3.}}
		Cho dãy $\{ u_n \}_{n=1}^{+\infty}$ được xác định bởi 
	\[ \left\{
	\begin{aligned}
		u_1 &= 2 \\
		u_{n+1} &= \dfrac{u_n^2 + 2003 u_n}{2004}, \ \forall n \in \mathbb{N}^*
	\end{aligned}
	\right. \]
	Đặt $v_n = \sum\limits_{i=1}^n \dfrac{u_i}{u_{i+1} - 1}$. Tính $\lim v_n$
	\end{BoxRed}
	\textbf{\textcolor{red}{LỜI GIẢI}} \\
    Bằng quy nạp dễ dàng có được $u_n > 1, \forall n \in \mathbb{N}^*$, từ đó
	$$ u_{n+1} = \dfrac{u_n (u_n - 1)}{2004} + u_n \implies u_{n+1} - u_n = \dfrac{u_n (u_n - 1)}{2004} \geq 0, \ \forall n \in \mathbb{N}^*$$
	$\implies \{u_n\}$ là một dãy tăng. Giả sử rằng $\{u_n\}$ bị chặn trên và hội tụ về $L > 2$. Cho qua giới hạn hai vế từ công thức truy hồi ta có: \\
	$ L = \dfrac{L^2 + 2003L}{2004} \implies L = 0 \lor L = 1$ (mâu thuẫn)\\
	Và từ công thức truy hồi, ta cũng có
	
	$ \dfrac{u_{n+1} - u_n}{u_n - 1} = \dfrac{u_n}{2004} \implies \dfrac{2004\left(u_{n+1} - u_n\right)}{\left(u_{n+1} - 1 \right)\left(u_n - 1 \right)} = \dfrac{u_n}{u_{n+1} - 1}$
	
	$\implies \dfrac{u_n}{u_{n+1} - 1} = 2024\left(\dfrac{1}{u_n - 1} - \dfrac{1}{u_{n+1} - 1} \right)$
	
	$\implies v_n = \sum\limits_{i=1}^n \dfrac{u_i}{u_{i+1} - 1} = 2004 \sum\limits_{i=1}^n \left(\dfrac{1}{u_i - 1} - \dfrac{1}{u_{i+1} - 1} \right) = 2004\left(\dfrac{1}{u_1 - 1} - \dfrac{1}{u_{n+1} - 1} \right)$
	
	$\implies \lim v_n = 2004$
    \newpage
    \begin{BoxRed}
	\textbf{\textcolor{red}{Bài toán 4.}}
		Cho dãy $\{ u_n \}_{n=1}^{+\infty}$ được xác định bởi
	\[ \left\{
	\begin{aligned}
		u_1 &= 1 \\
		u_{n+1} &= 1 + u_1 u_2\dots u_n, \ \forall n \in \mathbb{N}^*
	\end{aligned}
	\right. \]
	Đặt $v_n = \sum\limits_{k=1}^n \dfrac{1}{u_k}$. Tính $\lim v_n$
	\end{BoxRed}
	\textbf{\textcolor{red}{LỜI GIẢI}} \\
    Bằng quy nạp ta chứng minh được $u_n > 1, \ \forall n \geq 2$. \\ Từ công thức truy hồi ta có: $u_2 = 2 > 1 = u_1$\\
	$u_{n+1} - 1 = u_1u_2\dots u_n \implies u_n - 1 = u_1u_2\dots u_{n-1}, \ \forall n \geq 2$\\
	$\implies u_{n+1} = 1 + (u_n - 1)u_n \implies u_{n+1} - u_n = (u_n - 1)^2 > 0, \ \forall n \geq 2$\\
	$\implies \{u_n\}$ là một dãy tăng. Giả sử rằng $\{u_n\}$ bị chặn trên và hội tụ về $L > 1$. Cho qua giới hạn hai vế từ công thức truy hồi ta có \\
	$ L - L = (L - 1)^2 \implies L = 1$ (mâu thuẫn)\\
	$\implies \lim u_n = +\infty $\\
	Cũng từ công thức truy hồi ta có\\
    $u_{n+1} = 1 + (u_n - 1)u_n \implies u_{n+1} - 1 = (u_n - 1)u_n, \ \forall n \geq 2$

    $\implies \dfrac{1}{u_{n+1} - 1} = \dfrac{1}{u_n - 1} - \dfrac{1}{u_n} \implies \dfrac{1}{u_n} = \dfrac{1}{u_n - 1} - \dfrac{1}{u_{n+1} - 1}, \ \forall n \geq 2$

    $\implies v_n = \sum\limits_{k=1}^n \dfrac{1}{u_k} = u_1 + \sum\limits_{k=2}^n \left(\dfrac{1}{u_k - 1} - \dfrac{1}{u_{k+1} - 1} \right) = u_1 + \left(\dfrac{1}{u_2 - 1} - \dfrac{1}{u_{n+1} - 1} \right)$

    $\implies \lim v_n = 2$
    \newpage
    \begin{BoxRed}
		\textbf{\textcolor{red}{Bài toán 5.}}
		Cho dãy $\{ x_n \}_{n=1}^{+\infty}$ được xác định bởi
	\[ \left\{
	\begin{aligned}
		1 <& x_1 < 2 \\
		x_{n+1} &= 1 + x_n - \dfrac{x_n^2}{2}, \ \forall n \in \mathbb{N}^*
	\end{aligned}
	\right. \]
	Tính $\lim x_n$
	\end{BoxRed}
	\textbf{\textcolor{red}{LỜI GIẢI}} \\
    Ta sẽ chứng minh $1 < x_n < 2, \ \forall n \in \mathbb{N}^*$. Thật vậy, với $k=1$ thì ta có\\ $x_2 = 1 + x_1 - \dfrac{x_1^2}{2} \in (1, 2).$ Giả sử rằng điều này đúng đến $n=k,$ tức là \\ $1 < x_k < 2;$
    \\Ta cần chứng minh
    $ 1 < x_{k+1} < 2 \iff 1 < 1 + x_k - \dfrac{x_k^2}{2} < 2$\\
    BĐT cuối luôn đúng theo điều đã giả sử, nên theo nguyên lí quy nạp suy ra \\ $1 < x_n < 2,\ \forall n \in \mathbb{N}^*$\\
    $\implies x_n$ bị chặn. Giả sử rằng $\{x_n\}$ hội tụ về $1 < L < 2$. Cho qua giới hạn hai vế từ công thức truy hồi ta có
    $ L = 1 + L - \dfrac{x_n^2}{2} \implies L = \sqrt{2}$\\
    Ta sẽ chứng minh $x_n \rightarrow \sqrt{2}$. Thật vậy, xét hiệu

    $
    \begin{aligned}
    0 \leq \left|x_{n+1} - \sqrt{2}\right| &= \left|1 + x_n -\dfrac{x_n^2}{2} - \sqrt{2}\right| = \dfrac{1}{2}\left|x_n - \sqrt{2}\right|\left|2 - \sqrt{2}     - x_n\right| \\ 
    &=\dfrac{1}{2}\left|x_n - \sqrt{2}\right|\left|x_ 
    n - 2 + \sqrt{2}\right| < \dfrac{1}{\sqrt{2}}\left|x_n - \sqrt{2}\right| < \dots < \dfrac{1}{\sqrt{2}^n}\left|x_1 - \sqrt{2}\right|
    \end{aligned}$\\
    Mặt khác $\lim \dfrac{1}{\sqrt{2}^n}\left|x_1 - \sqrt{2}\right| = 0$ (vì $x_1$ là một hằng số hữu hạn) \\
    Theo BĐT trên và nguyên lí kẹp suy ra $ \ x_n \rightarrow \sqrt{2} $ (ĐPCM)
    \newpage
    \begin{BoxRed}
	\textbf{\textcolor{red}{Bài toán 6.}} Giải hệ phương trình 
		\[\left\{
        \begin{gathered}
            x - y -3 = 2\sqrt{x^2 + 3} - 2\sqrt{y^2 + 6y + 12} \\
            (2x -5)\sqrt{x + 2} + (2y + 7)\sqrt{1 - y} + \sqrt{xy + 2x - y -2} = \left|y + 2\right| + 4x - 4
        \end{gathered}
        \right.\]
	\end{BoxRed}
	\textbf{\textcolor{red}{LỜI GIẢI}} \\
    ĐKXĐ: $x \geq -2, y \leq 1.$
	Gọi các phương trình của hệ lần lượt là $(1), (2).$
    
    Ta có $(1) \iff 2\sqrt{x^2 + 3} - x = 2\sqrt{y^2 + 6y + 12} - (y + 3)\\$
    Xét hàm số $f(t) = 2\sqrt{t^2 + 3} - t \implies f'(t) = \dfrac{t}{\sqrt{t^2 + 3}} - 1.$\\
    Mặt khác $\sqrt{t^2 + 3} > |t| \geq t \implies \dfrac{t}{\sqrt{t^2 + 3}} < 1 \implies f'(t) < 0, \ \forall t \in \mathbb{R}$\\
    $\implies f(x) = f(y+3) \iff x = y + 3,\ y \geq -5$ Thay vào $(2)$ ta được
    
    $(2y + 1)\sqrt{y+5} + (2y+7)\sqrt{1-y} + \sqrt{(y+2)^2} = |y+2| + 4y + 8$
    
    $\iff (2y + 1)\sqrt{y+5} + (2y+7)\sqrt{1-y} = 4y + 8$
    
    Đặt $a = \sqrt{y + 5} \geq 0,\ b = \sqrt{1 - y} \geq0$, khi đó ta có:
    
    $\left\{
        \begin{gathered}
        \begin{aligned}
            a^2 + b^2 &= 6 &\text{(1)}\\
            (2a^2 - 9)a + (9 - 2b^2)b &= 2(a^2 - b^2) &\text{(2)}
        \end{aligned}
        \end{gathered}
    \right.$
    
    $\begin{aligned}
        (2) &\iff (a-b)(2a^2 + 2ab + 2b^2 - 9 - 2a - 2b)=0 \\
            &\iff a - b = 0 \lor 2a^2 + 2ab + 2b^2 - 9 - 2a - 2b = 0\\
            &\iff a = b \lor (a+b+1)(a+b-3) = 0\\
            &\iff a = b \lor a+b = 3\\
            &\iff \sqrt{y+5} = \sqrt{1-y} \lor \sqrt{y+5}+\sqrt{1-y}=3\\
            &\iff y = -2 \lor y = -\dfrac{4+3\sqrt{3}}{2} \lor y = \dfrac{3\sqrt{3}-4}{2}\\
            &\iff x = 1 \lor x = \dfrac{2 - 3\sqrt{3}}{2} \lor x = \dfrac{3\sqrt{3} + 2}{2}
    \end{aligned}$\\
    Vậy hệ phương trình có các nghiệm là:
    \[(x; y) \in \left\{ \left(1 ; -2\right); \left(\dfrac{2 - 3\sqrt{3}}{2}; -\dfrac{4+3\sqrt{3}}{2}\right); \left(\dfrac{3\sqrt{3} + 2}{2}; \dfrac{3\sqrt{3}-4}{2}\right) \right\}\]
    \newpage
    \begin{BoxRed}
	\textbf{\textcolor{red}{Bài toán 7.}} 
          Cho tập hợp $\Omega$ là họ các tập con của tập $ \left\{m \in \mathbb{N^*} \ | \ m \leq 2n\right\},$ hãy tìm số lượng tập con có thể có sao cho không tồn tại cặp số $(a; b)$ mà $|a - b| = n$ trong các tập con nói trên.

	\end{BoxRed}
	\textbf{\textcolor{red}{LỜI GIẢI}} \\
    Gọi $A_i$ là tập con của $\Omega$ tương ứng có chứa cặp số $(n+i; i)$, trong đó $i = \overline{1, n}$.
    Khi đó các tập vừa được định nghĩa ở trên là các tập không thỏa yêu cầu bài toán. Như vậy ta có $n$ cặp số $(a; b) \in \Omega$ mà $|a - b| = n$. Ta có:\\
    $|\Omega| = 2^{2 n} = 4^n = C_n^0 4^n$ \\   
    $\bigcup\limits_{i=1}^n A_i = \sum\limits_{i=1}^n |A_i| - \sum\limits_{1\leq i < j \leq n}|A_i \cap A_j| + \sum\limits_{1\leq i < j < k \leq n} |A_i \cap A_j \cap A_k| - \dots$ \\
    Mặt khác vai trò của $A_1, A_2, \dots$ là như nhau nên ta có:\\
    $\begin{aligned}
       & \ \ \ \sum \limits_{i=1}^n |A_i| = C_n^1 |A_1| = C_n^1 2^{2n-2} = C_n^1 4^{n-1}\\
        &\sum \limits_{1\leq i < j \leq n}|A_i \cap A_j| = C_n^2 |A_1 \cap A_2| = C_n^2 2^{2n-4} = C_n^2 4^{n-2}   
    \end{aligned}$
   
    $\dots$
    $\implies \bigcup\limits_{i=1}^n A_i = C_n^1 4^{n-1} - C_n^2 4^{n-2} + \dots +(-1)^n C_n^n 4^0$\\
    $\implies |\Omega| - \bigcup\limits_{i=1}^n A_i = C_n^0 4^n - C_n^1 4^{n-1} + C_n^2 4^{n-2} - \dots +(-1)^n C_n^n 4^0 = 3^n$\\
    Vậy, theo nguyên lý bù trừ suy ra số tập con cần tìm là $3^n$
    \newpage
    \begin{BoxRed}
	\textbf{\textcolor{red}{Bài toán 8.}} 
          Tìm tất cả các hàm số $f: \mathbb{Z} \rightarrow \mathbb{R}$ sao cho 
          \[\left\{
        \begin{gathered}
            f(1) = \dfrac{5}{4}\quad\quad\quad\quad\quad\quad\quad\quad\quad\quad\quad\quad\quad\quad\quad\quad \ \ \ (1)\\ 
            f(m + n) + f(m - n) = 2f(m)f(n), \forall m, b \in \mathbb{Z} \quad (2)
        \end{gathered}    
        \right.\]

	\end{BoxRed}
	\textbf{\textcolor{red}{LỜI GIẢI}} \\
    Giả sử tồn tại hàm số thỏa mãn yêu cầu bài toán.
    Kí hiệu $P(a, b)$ là phép thế cặp số $(a; b)$ vào $(2)$.\\
    $P(n, 0) \implies 2f(n) = 2f(n)f(0) \implies f(0) = 1 \lor f(n) = 0$.\\
    Dễ thấy $f(n) = 0, \ \forall n \in \mathbb{Z}$ là điều vô lí. 
    Vậy $f(0) = 0 \quad (3)$\\
    $\text{Từ $(3)$ và $P(0; n)$}\\ 
    \implies f(n) + f(-n) = 2f(n) \implies f(n) = f(-n), \ \forall n \in \mathbb{Z}\\
    \implies f$ là hàm chẵn\\
    Từ $(1)$ và $P(n; 1)$\\
    $\implies f(n+1) + f(n-1) = \dfrac{5}{2}f(n) \\ \implies f(n+1) = \dfrac{5}{2}f(n) - f(n-1), \ \forall n \in \mathbb{Z} \quad (4)$\\
    Ta sẽ chứng minh mệnh đề: $f(n) = \dfrac{1}{2}\left(2^n + \dfrac{1}{2^n}\right), \ \forall n \in \mathbb{N}$.\\
    Thật vậy, với $n = 0$ mệnh đề đúng. \\
    Giả sử mệnh đề đúng tới $n$ nào đó, tức là $f(n) = \dfrac{1}{2}\left(2^n + \dfrac{1}{2^n}\right)$\\
    Từ $(4)$ ta có:\\
    $\begin{aligned}
        f(n+1) &= \dfrac{5}{2}f(n) - f(n-1) = \dfrac{5}{4}\left(2^n + \dfrac{1}{2^n}\right) - \dfrac{1}{2}\left(2^{n-1} + \dfrac{1}{2^{n-1}}\right)\\
               &= 2^n \left(\dfrac{5}{4} - \dfrac{1}{4}\right) + \dfrac{1}{2^n} \left(\dfrac{5}{4} - 1\right) = \dfrac{1}{2} \left(2^{n+1} - \dfrac{1}{2^{n+1}}\right)                
    \end{aligned}$\\
    Vậy theo nguyên lý quy nạp ta có: $f(n) = \dfrac{1}{2}\left(2^n + 2^{-n}\right), \ \forall n \in \mathbb{Z}$.\\
    Thử lại ta thấy thỏa mãn.
    \newpage
    \begin{BoxRed}
	\textbf{\textcolor{red}{Bài toán 9.}} 
         Cho dãy $(a_n)$ được xác định bởi
	\[ \left\{
	\begin{aligned}
		&a_1 = 1 \\
		&a_{n+1} = \sqrt{a_n + n + 1}, \ \forall n \in \mathbb{N}^*
	\end{aligned}
	\right. \]
    Tính $\lim \left(a_n - \sqrt{n}\right)$.
	\end{BoxRed}
	\textbf{\textcolor{red}{LỜI GIẢI}} \\
    Trước hết ta sẽ chứng minh mệnh đề $P(n)$: "$a_n \leq \sqrt{n-1} + 1$", $\forall n \in \mathbb{N^*}$. Thật vậy, ta thấy $P(1)$ đúng. Giả sử mệnh đề đúng tới $n$ bất kì, tức là $a_n \leq \sqrt{n-1} + 1$, ta có:\\
    $a_{n+1} = \sqrt{a_n + n + 1} \leq \sqrt{\sqrt{n-1}+n+2}$. Ta sẽ chứng minh \\
    $\sqrt{\sqrt{n-1}+n+2} \leq \sqrt{n} + 1 \iff \sqrt{n-1}+1\leq 2\sqrt{n}$ (đúng với mọi $n\geq 1$).\\
    Vậy $P(n+1)$ đúng. Do đó theo quy nạp, phép chứng minh hoàn tất.\\
    Tương tự ta cũng chứng minh được $a_n \geq \sqrt{n}, \ \forall n \in \mathbb{N^*}$\\
    $\implies 0 \leq a_n - \sqrt{n} \leq \sqrt{n-1} + 1 - \sqrt{n}, \ \forall n \in \mathbb{N^*}\\
    \implies 0 \leq \dfrac{a_n - \sqrt{n}}{\sqrt{n}} \leq \dfrac{\sqrt{n-1} + 1 - \sqrt{n}}{\sqrt{n}} = \dfrac{1}{\sqrt{n}} + \dfrac{\sqrt{n-1}}{\sqrt{n}} - 1$\\
    Dễ thấy $\lim 0 = \lim \dfrac{1}{\sqrt{n}} + \dfrac{\sqrt{n-1}}{\sqrt{n}} - 1 = 0$ nên theo nguyên lý kẹp suy ra $\lim \dfrac{a_n - \sqrt{n}}{\sqrt{n}} = 0$.\\
    Đặt $b_n = a_n - \sqrt{n}$ khi đó: $\lim \dfrac{b_n}{\sqrt{n}} =0, a_n = b_n + \sqrt{n}$. Từ đây ta có:\\
    $b_{n+1} = \sqrt{b_n + \sqrt{n} + n + 1} -\sqrt{n+1}, \ \forall n \in \mathbb{N}^*$
    
	$\begin{aligned}
    \implies \lim b_n &= \lim\left(\sqrt{b_n + \sqrt{n} + n + 1} -\sqrt{n+1}\right) = \lim \dfrac{b_n + \sqrt{n}}{\sqrt{b_n + \sqrt{n} + n + 1} +\sqrt{n+1}}\\
                      &= \lim \dfrac{\dfrac{b_n}{\sqrt{n}}+1}{\sqrt{\dfrac{b_n}{n}+\dfrac{1}{\sqrt{n}}+1 + \dfrac{1}{n}}+ \sqrt{1 + \dfrac{1}{n}}} = \dfrac{1}{2}
	\end{aligned}$
    
    Vậy $\lim \left(a_n - \sqrt{n}\right) = \dfrac{1}{2}$
    \newpage
    \begin{BoxRed}
	\textbf{\textcolor{red}{Bài toán 10 - Bổ đề 1.}} 
         Cho số thực $q \in (0; 1)$ và hai dãy không âm $(a_n), \ (b_n)$ thoả mãn
         $\lim b_n = 0, \ a_{n+1} \leq q a_n + b_n, \ \forall n \in \mathbb{N}$.\\ 
    Chứng minh rằng $\lim a_n = 0$
	\end{BoxRed} 
	\textbf{\textcolor{red}{LỜI GIẢI}} \\
    
\end{document}
